% DeepPass21: detailed projection/resolvent identities.

对自伴算符 $A=QHQ$，谱定理给出
\begin{equation}
 \boxed{
 A=\int_{\sigma(A)}\lambda\,\dd E_A(\lambda),
 \qquad
 (z-A)^{-1}=\int_{\sigma(A)}\frac{\dd E_A(\lambda)}{z-\lambda}.}
 \label{eq:resolvent-spectral-resolution-dp68}
\end{equation}
由此可把预解式的算符范数直接写成到谱的距离：
\begin{equation}
\boxed{
\|(z-QHQ)^{-1}\|
=\frac1{\operatorname{dist}[z,\sigma(QHQ)]}}
\qquad
(z\notin\sigma(QHQ),\ QHQ=QHQ^\dagger).
\label{eq:resolvent-norm-spectral-dp21}
\end{equation}

若 $A$ 可逆，则
\begin{equation}
 \boxed{A-B=(I-BA^{-1})A.}
 \label{eq:neumann-factorization-dp68}
\end{equation}
当 $\|BA^{-1}\|<1$ 时，几何级数在算符范数下收敛，得到
\begin{equation}
\boxed{
(A-B)^{-1}=A^{-1}\sum_{n=0}^{\infty}(BA^{-1})^n}
\qquad
\text{if }\|BA^{-1}\|<1.
\label{eq:operator-neumann-resolvent-dp21}
\end{equation}

若 $P$ 中只保留一个非简并裸态 $|p\rangle$，且 $P$--$Q$ 混合足够弱，精确 Feshbach 自能的最低非零静态极限为
\begin{equation}
 \boxed{
 \Delta E_p^{(2)}
 =\sum_{q\in Q}
 \frac{|H_{pq}|^2}{E_p-E_q}.}
 \label{eq:feshbach-second-order-static-dp68}
\end{equation}
这就是普通非简并二阶微扰论在投影预解式中的来源。
